\chapter{Marco teórico referencial de la investigación}\label{chap:1}

El presente capítulo tiene como objetivo exponer los fundamentos teóricos que sustentan la investigación, abordando los conceptos, tecnologías y enfoques necesarios para el diseño de un sistema basado en el Internet de las Cosas (IoT) aplicado a la agricultura. En particular, se analizan los principios del IoT en entornos agrícolas, los sistemas de monitoreo para variables ambientales y las tecnologías de transmisión de datos en tiempo real, su aplicación en contextos con recursos limitados. Asimismo, se revisan soluciones existentes con el fin de identificar enfoques, herramientas y limitaciones relevantes para el desarrollo de la propuesta planteada en esta investigación.

% ─────────────────────────────────────────────
\section{Situación problemática}\label{sec:situacion}
% ─────────────────────────────────────────────

La agricultura cubana enfrenta una crisis de carácter multiplicativo en la que la escasez de materias primas, energía, agua y suelo se retroalimenta: cada factor degradado amplifica el impacto de los demás \citep{cigg2025}. En el plano hídrico, el país registra sequías prolongadas, acumulados de lluvia por debajo de los valores medios históricos y un patrón errático de precipitaciones ---escasez prolongada interrumpida por tormentas torrenciales que dañan físicamente cultivos de plátano, maíz y frutales, además de provocar erosión severa. El cambio climático agrava estas condiciones mediante olas de calor, huracanes e inundaciones que alteran la planificación agrícola.

Paradójicamente, la infraestructura de riego ---el recurso más necesario para combatir la sequía--- se encuentra colapsada. Solo una fracción menor de los cultivos dispone de sistemas de riego operativos, predominando tecnologías obsoletas como el riego por aspersión, gravedad, aniego y tracción animal, lo que se traduce en bajo rendimiento por hectárea \citep{cornelltraccion}. A esto se suma la ausencia de precisión agrícola: siembra a voleo, fertilización uniforme sin análisis de suelo y gestión del riego ``por calendario'' en lugar de por necesidad real del cultivo.

La escasez de combustible agrava aún más el problema, pues sin electricidad no es posible bombear agua de embalses ni pozos hacia los canales de riego, y la maquinaria de bombeo permanece paralizada por falta de piezas \citep{alnap2024}. Adicionalmente, la sobreexplotación de suelos marginales sin fertilizantes ni drenaje adecuado ha avanzado la salinización, tornando ineficaz incluso el agua disponible. Ante este panorama estructural ---donde la crisis hídrica, el colapso de infraestructura y las restricciones externas se refuerzan mutuamente--- resulta imperativo explorar enfoques tecnológicos que, por su costo, apertura y autonomía energética, sean viables dentro de estas condiciones y no a pesar de ellas. La sección siguiente revisa los antecedentes y las tecnologías que han evolucionado en esa dirección.

Subyacente a todas estas carencias se encuentra el bloqueo económico, comercial y financiero impuesto por los Estados Unidos, que en su configuración actual comprende 246 medidas coercitivas de carácter extraterritorial \citep{cubabloqueo2024}. Estas medidas restringen el acceso de Cuba a mercados internacionales de tecnología, equipos agropecuarios, repuestos, combustibles y divisas, impidiendo que instituciones como el MINAG o el INRH adquieran insumos básicos directamente de proveedores en países terceros bajo amenaza de sanciones secundarias. El resultado práctico es que soluciones tecnológicas disponibles comercialmente a nivel global ---sistemas de riego por goteo, sensores industriales, bombas eficientes--- resultan inaccesibles no por 
desconocimiento técnico ni por ausencia de voluntad institucional, sino por una restricción de carácter político impuesta desde el exterior. Este contexto convierte la selección de tecnologías abiertas, de componentes ampliamente distribuidos y de arquitecturas que minimicen la dependencia de importaciones especializadas en una condición de diseño no negociable para cualquier propuesta viable en el país.

% ─────────────────────────────────────────────
\section{Antecedentes y tecnologías actuales}\label{sec:antecedentes}
% ─────────────────────────────────────────────

Antes de que existiera el término ``Internet de las Cosas'', ya se empleaban sistemas de telemetría para monitorear variables agrícolas a distancia. Estos sistemas se apoyaban en radiofrecuencia o líneas telefónicas y eran ofrecidos por empresas como Rain Bird o Toro mediante controladores de riego con comunicación propietaria de alto costo. El concepto de \emph{redes de sensores inalámbricos} (WSN, por sus siglas en inglés), considerado el precursor directo del IoT, fue impulsado por proyectos académicos y de defensa como LOFAR, orientado a medir temperatura y humedad del suelo y transmitir los datos vía radio hacia una estación base.

Con la masificación de Internet a finales de los años noventa, los primeros objetos conectados surgieron en entornos de automatización industrial bajo protocolos como Modbus y OPC. No fue sino hasta 1999 cuando Kevin Ashton acuñó el término \emph{Internet of Things} para describir la visión de conectar objetos físicos a la red global mediante identificación por radiofrecuencia (RFID). Desde entonces, el ecosistema IoT ha experimentado una aceleración notable: el reporte del IEEE de marzo de 2014 lo define como una red de elementos físicos, cada uno dotado con sensores, interconectados a través de Internet \citep{ieeereport2014}.

En agricultura, esta evolución se tradujo primero en estaciones meteorológicas autónomas y sistemas SCADA para grandes explotaciones. Posteriormente, la reducción de costos de microcontroladores, módulos de comunicación inalámbrica y sensores ambientales democratizó el acceso a estas soluciones. Plataformas como Arduino, aparecida en 2005, y servicios en la nube de bajo costo permitieron que investigadores y agricultores de países en desarrollo comenzaran a desarrollar sistemas propios, adaptados a sus restricciones locales \citep{arduinoopensource}. En la actualidad, el mercado ofrece una amplia gama de opciones ---desde microcontroladores de cuatro dólares hasta gateways industriales con 5G integrado--- cuyo análisis resulta indispensable para seleccionar la arquitectura más adecuada al contexto cubano. Los fundamentos de estas tecnologías se examinan en la sección siguiente.

% ─────────────────────────────────────────────
\section{Fundamentos teóricos}\label{sec:fundamentos}
% ─────────────────────────────────────────────

% ── 1.3.1 ──────────────────────────────────────
\subsection{Internet de las Cosas (IoT) aplicado a la agricultura}\label{subsec:iot-agricultura}

El Internet de las Cosas representa un paradigma tecnológico que integra dispositivos físicos, conectividad y procesamiento de datos para generar valor en sectores productivos \citep{ieeereport2014}. En el ámbito agrícola, el IoT ha emergido como una herramienta poderosa para optimizar el uso de recursos y mejorar la productividad \citep{quiroz2024transformacion}, pues permite la recopilación de datos ambientales en tiempo real y facilita la toma de decisiones informadas para la gestión de cultivos.

El ecosistema IoT se estructura en múltiples capas que deben integrarse de forma coherente. La \textbf{capa de hardware} comprende los microcontroladores y plataformas SoC encargados de la adquisición y el procesamiento inicial de datos. Entre las opciones de mayor relevancia para aplicaciones agrícolas de bajo costo se encuentran:

\begin{itemize}
	\item \textbf{Arduino UNO}: plataforma electrónica de código abierto basada en hardware y software de fácil uso \citep{arduinoopensource}. Dispone de catorce pines de entrada/salida digital, seis entradas analógicas, conexión USB, cristal de cuarzo de 16~MHz y capacidad de programación en circuito. Su entorno de desarrollo integrado es compatible con Windows, macOS y Linux, e incluye numerosas bibliotecas que reducen la complejidad de programación. Su principal ventaja en contextos con restricciones presupuestarias radica en que las placas son considerablemente más económicas que otras plataformas, lo que elimina la preocupación por costos de reemplazo ante daños accidentales \citep{arduinoopensource}.
	\item \textbf{ESP32}: integra conectividad WiFi y Bluetooth de modo dual-core, con un costo aproximado de tres a cinco dólares, lo que lo ha convertido en una opción popular para proyectos IoT \citep{quiroz2024transformacion}.
	\item \textbf{Raspberry Pi Pico}: equipado con el microcontrolador RP2040 y procesador dual-core ARM Cortex-M0+, a un costo de aproximadamente cuatro dólares, con capacidades de programación de entrada/salida flexible.
\end{itemize}

La \textbf{capa de conectividad} determina el alcance, la velocidad de transmisión, el consumo energético y la viabilidad económica del sistema. Las tecnologías de mayor pertinencia para agricultura en zonas rurales con infraestructura limitada son:

\begin{itemize}
	\item \textbf{LoRa/LoRaWAN}: alcances superiores a 15~km con consumo ultra-bajo, velocidades entre 0,3 y 50~kbps, particularmente útil en agricultura inteligente y seguimiento en zonas rurales.
	\item \textbf{WiFi (802.11 b/g/n)}: alta velocidad y amplia disponibilidad de infraestructura, adecuado cuando existe cobertura local.
	\item \textbf{NB-IoT}: opera sobre redes celulares licenciadas con penetración profunda en interiores y autonomía de batería superior a diez años, orientado a medidores inteligentes y sensores subterráneos.
\end{itemize}

La \textbf{capa de software y protocolos} define cómo los datos fluyen desde los dispositivos hacia los sistemas de procesamiento y visualización. El protocolo MQTT (\emph{Message Queue Telemetry Transport}) destaca como estándar de facto en IoT agrícola: es un protocolo ligero basado en publicación-suscripción que opera sobre TCP/IP, diseñado específicamente para dispositivos con recursos limitados y redes de alta latencia \citep{tian2024novel}. Su modelo utiliza un broker central que gestiona y distribuye los mensajes a los suscriptores correspondientes, lo que simplifica la arquitectura de comunicación.

La arquitectura de referencia para sistemas IoT agrícolas típicamente sigue un modelo de cuatro capas: sensores, comunicación, internet y aplicación \citep{shahroorei2024smart}. En el contexto de este trabajo, dicho modelo se implementa con sensores de bajo costo conectados a un microcontrolador Arduino UNO, que actúa como nodo de adquisición y control local, transmitiendo datos mediante MQTT hacia un gateway de bajo consumo para su visualización y análisis. El análisis de los sensores que componen la capa de percepción se aborda en la subsección siguiente.

% ── 1.3.2 ──────────────────────────────────────
\subsection{Sistemas de monitoreo de variables ambientales}\label{subsec:monitoreo}

El monitoreo preciso de variables ambientales constituye el fundamento operativo de cualquier sistema de riego inteligente. Para la gestión hídrica en agricultura, las variables críticas comprenden la temperatura y humedad del aire, la humedad del suelo, la presión atmosférica y, en algunos contextos, la radiación solar y la velocidad del viento. La selección de sensores adecuados debe equilibrar precisión, durabilidad, consumo energético y costo de adquisición, criterio determinante en el contexto cubano.

El \textbf{sensor DHT22} es uno de los más empleados en sistemas IoT agrícolas de bajo costo. Mide temperatura en un rango de -40~°C a +80~°C con precisión de ±0,5~°C y humedad relativa entre 0\% y 100\% con precisión de ±2\%. Su interfaz digital de un solo cable simplifica la integración con Arduino. Estudios recientes han validado su efectividad en el monitoreo en tiempo real de cultivos de arroz, combinado con la plataforma ThingSpeak y protocolos HTTP y MQTT \citep{lozano2024desarrollo}. El \textbf{BMP180} complementa estas mediciones aportando presión atmosférica y temperatura barométrica mediante interfaz I$^2$C, útil para estimar altitud y correlacionar condiciones climáticas.

Para la medición de humedad del suelo existen dos enfoques principales. Los \textbf{sensores resistivos} son los de menor costo y mayor disponibilidad, pero presentan deriva a largo plazo por corrosión de los electrodos. Los \textbf{sensores capacitivos} ofrecen mayor durabilidad y estabilidad, siendo preferibles para despliegues permanentes. Ambos tipos entregan una señal analógica proporcional al contenido de agua en el suelo, que el microcontrolador convierte mediante su conversor analógico-digital.

La \textbf{gestión de datos} generados por estos sensores requiere soluciones de almacenamiento adaptadas a la naturaleza temporal de la información. Las \textbf{bases de datos de series temporales} han emergido como la solución predominante para aplicaciones IoT \citep{shahroorei2024smart}. Para dispositivos con recursos severamente limitados, \textbf{SQLite} ---con un tamaño aproximado de 600~kB--- proporciona SQL completo sin configuración previa, operaciones rápidas de lectura y escritura, y cumplimiento de propiedades ACID que garantizan consistencia incluso ante fallos de energía \citep{shahroorei2024smart}, características especialmente valiosas en un entorno con cortes eléctricos frecuentes como el cubano. Para soluciones de mayor escala, \textbf{InfluxDB} se ha posicionado como líder en el ámbito IoT, ofreciendo un lenguaje de consultas específico y alta compresión de datos bajo licencia MIT \citep{shahroorei2024smart}.

La precisión y confiabilidad de estas mediciones depende, en última instancia, de la capacidad del sistema para transmitir los datos de forma robusta y oportuna. Los mecanismos de transmisión disponibles para entornos con recursos limitados se analizan a continuación.

% ── 1.3.3 ──────────────────────────────────────
\subsection{Transmisión de datos en tiempo real en entornos con recursos limitados}\label{subsec:transmision}

La transmisión de datos en tiempo real en contextos de recursos limitados ---caracterizados por ancho de banda reducido, alta latencia, interrupciones frecuentes y restricciones energéticas--- requiere una selección cuidadosa de protocolos y tecnologías de red. El diseño de esta capa determina en gran medida la viabilidad operativa del sistema completo.

El protocolo \textbf{MQTT} es la opción predominante para este tipo de escenarios. Opera sobre TCP/IP con una sobrecarga de encabezado de apenas dos bytes, lo que lo hace significativamente más eficiente que HTTP en condiciones de baja conectividad \citep{tian2024novel}. Su arquitectura de publicación-suscripción desacopla productores y consumidores de datos, permitiendo que el sistema continúe funcionando localmente aunque la conexión a Internet se interrumpa, siempre que el broker opere en la red local. Como alternativa ultraligera, \textbf{CoAP} (\emph{Constrained Application Protocol}) implementa un modelo RESTful sobre UDP/IP, proporcionando comunicación de baja sobrecarga similar a HTTP pero optimizada para redes restringidas \citep{rfc7252}.

En cuanto a la \textbf{tecnología de red física}, la elección depende de la distancia entre nodos y la infraestructura disponible. Para instalaciones donde los sensores se ubican en un radio menor a 100~m de un gateway con acceso a corriente eléctrica, WiFi resulta la opción más práctica por su disponibilidad y facilidad de configuración. Para parcelas extensas o zonas sin infraestructura de red, \textbf{LoRaWAN} ofrece alcances de hasta 15~km con consumo de apenas microamperios en modo de bajo consumo, siendo especialmente pertinente para el contexto rural cubano.

Un aspecto crítico en entornos con cortes eléctricos frecuentes es la capacidad de \textbf{operación autónoma local}. La arquitectura edge-first, en la que el microcontrolador toma decisiones de riego de forma autónoma con base en los datos de los sensores locales ---sin depender de conectividad a la nube--- garantiza la continuidad operativa ante fallos de red o energía. Los datos se almacenan localmente en SQLite y se sincronizan con el servidor central cuando la conectividad se restablece. Este patrón de diseño, conocido como \emph{store-and-forward}, resulta fundamental para la resiliencia del sistema en el contexto cubano.

El software de código abierto refuerza la viabilidad de estas soluciones en países en desarrollo, proporcionando alternativas escalables y económicas sin las barreras financieras del software propietario \citep{livingopensource2024impact}. La combinación de protocolos abiertos, almacenamiento embebido y hardware de código abierto como Arduino configura un stack tecnológico sostenible. La revisión de implementaciones concretas que integran estos elementos se presenta en la sección siguiente.

% ─────────────────────────────────────────────
\section{Soluciones similares}\label{sec:similares}
% ─────────────────────────────────────────────

El análisis de soluciones existentes permite identificar enfoques validados, limitaciones recurrentes y oportunidades de mejora aplicables al diseño del sistema propuesto. En la literatura reciente se distinguen tres categorías principales de implementaciones IoT para riego agrícola.

\textbf{Sistemas comerciales de precisión agrícola}, como los ofrecidos por empresas Trimble, John Deere o Netafim, integran sensores de suelo, estaciones meteorológicas y actuadores de riego con plataformas de analítica en la nube. Si bien demuestran alta efectividad en la reducción del consumo de agua ---entre 20\% y 40\% respecto al riego convencional---, presentan costos de adquisición e instalación que los hacen inaccesibles para el contexto cubano, además de depender de conectividad celular estable y licencias de software propietario.

\textbf{Sistemas académicos de bajo costo} orientados a países en desarrollo han demostrado la viabilidad de aproximaciones más accesibles. Un estudio relevante demostró que el uso de sensores como el DHT22, el BMP-180 y sensores de humedad de suelo, combinados con la plataforma ThingSpeak y protocolos HTTP y MQTT, resultó efectivo para el monitoreo en tiempo real de cultivos de arroz, proporcionando información suficiente para ajustar las prácticas agrícolas según las condiciones requeridas \citep{lozano2024desarrollo}. Otro trabajo confirmó que la implementación de modelos IoT con sensores de bajo costo produce mejoras significativas en eficiencia y rentabilidad, concluyendo que el monitoreo en tiempo real se convierte en herramienta esencial para la adaptación de cultivos ante el cambio climático \citep{lozano2024desarrollo}.

\textbf{Sistemas con Arduino y conectividad local} representan la categoría más próxima a los requerimientos del presente trabajo. Diversas implementaciones han explorado el uso de Arduino UNO con módulos WiFi ESP8266 o NRF24L01 para transmitir datos de sensores de humedad y temperatura hacia servidores locales. No obstante, la mayoría de estos trabajos no abordan la resiliencia ante cortes eléctricos ni la operación autónoma prolongada sin intervención humana, aspectos críticos para el contexto cubano.

El análisis comparativo de estas soluciones revela tres limitaciones recurrentes: la dependencia de infraestructura de conectividad estable, la ausencia de mecanismos de operación autónoma ante fallos energéticos, y la falta de adaptación a restricciones de importación y disponibilidad de componentes. Estas brechas fundamentan la propuesta del sistema FoT (\emph{Farm of Things}), cuyas tecnologías seleccionadas se justifican en las conclusiones de este capítulo.

% ─────────────────────────────────────────────
\section{Conclusiones parciales}\label{sec:conclusiones-cap1}
% ─────────────────────────────────────────────

Una vez concluido el análisis del marco teórico referencial, se arriban a las siguientes conclusiones:

\begin{enumerate}
	\setlength\itemsep{0em}
	\item La solución propuesta debe cumplir los siguientes requisitos principales: bajo costo de adquisición y reemplazo de componentes, capacidad de operación autónoma local ante cortes eléctricos mediante el patrón \emph{store-and-forward}, consumo energético mínimo compatible con fuentes alternativas como baterías o paneles solares, y empleo exclusivo de software y hardware de código abierto que no dependan de licencias comerciales ni importaciones especializadas.
	\item La revisión de soluciones similares evidencia que ninguna de las implementaciones analizadas combina simultáneamente los requisitos de bajo costo, resiliencia ante interrupciones eléctricas frecuentes, operación autónoma prolongada y adaptación a las restricciones de importación propias del contexto cubano, lo que justifica el desarrollo de una solución propia.
	\item Las tecnologías seleccionadas para el sistema FoT son: \textbf{Arduino UNO} como plataforma microcontroladora, por su costo reducido, ecosistema amplio y disponibilidad en mercados alternativos; los sensores \textbf{DHT22} y \textbf{capacitivo de humedad de suelo}, por su precisión suficiente, bajo costo y amplia documentación; el protocolo \textbf{MQTT} sobre red WiFi local, por su ligereza y capacidad de operación con broker embebido sin dependencia de Internet; y \textbf{SQLite} como base de datos embebida, por sus propiedades ACID que garantizan consistencia ante fallos de energía con mínimos requerimientos de recursos.
\end{enumerate}

\pagebreak